\documentclass[a4paper,12pt]{article}
\usepackage[T2A]{fontenc}
\usepackage[utf8]{inputenc}
\usepackage[russian]{babel}
\usepackage{amsmath,amssymb,mathtools}
\usepackage{helvet}
\renewcommand{\familydefault}{\sfdefault}
\usepackage{geometry}
\geometry{left=20mm,right=19mm,top=21mm,bottom=22mm}
\usepackage{microtype}
\usepackage{booktabs,tabularx,array}
\usepackage{enumitem}
\usepackage{xcolor}
\usepackage{tikz}
\usepackage{hyperref}
\usepackage{fancyhdr}
\usepackage{parskip}
\usetikzlibrary{calc}
\definecolor{accent}{HTML}{007F85}
\definecolor{deepaccent}{HTML}{00565C}
\definecolor{bodyink}{HTML}{262B2F}
\definecolor{muted}{HTML}{56646A}
\definecolor{errorred}{HTML}{A23A39}
\definecolor{ideablue}{HTML}{245E8B}
\hypersetup{unicode=true,colorlinks=true,linkcolor=deepaccent,pdfauthor={\foreignlanguage{russian}{Лёвин Артём Александрович}},pdftitle={\foreignlanguage{russian}{Ревью домашней работы — Володя Хачатурян — 16.11.26}}}
\setlength{\parindent}{0pt}
\setlength{\parskip}{5pt plus 1pt minus 1pt}
\setlist[itemize]{leftmargin=1.5em,itemsep=2pt,topsep=3pt}
\setlist[enumerate]{leftmargin=1.5em,itemsep=2pt,topsep=3pt}
\pagestyle{fancy}
\fancyhf{}
\fancyhead[L]{\small\color{muted}Физика · Динамика и трение}
\fancyhead[R]{\small\color{muted}16 ноября 2026 г.}
\fancyfoot[C]{\small\color{muted}\thepage}
\renewcommand{\headrulewidth}{0.25pt}
\setlength{\headheight}{15pt}
\newcommand{\reviewtitle}[1]{\noindent{\Large\bfseries\color{deepaccent}Задание №#1}\par\vspace{2mm}\noindent\rule{\linewidth}{0.35pt}\vspace{2mm}}
\newcommand{\subblock}[1]{\par\medskip\textbf{\color{deepaccent}#1}\par\vspace{-2pt}}
\newcommand{\errorlabel}{\textcolor{errorred}{\textbf{Ошибка:}}}
\newcommand{\keyidea}{\textcolor{ideablue}{\textbf{Ключевая идея:}}}
\newcommand{\resultlabel}{\textcolor{deepaccent}{\textbf{Итог:}}}
\begin{document}

% Титульный лист
\thispagestyle{empty}
\begin{center}
\vspace*{29mm}
{\color{deepaccent}\large ФИЗИКА}\par
\vspace{17mm}
{\fontsize{25}{30}\selectfont Ревью домашней работы}\par
\vspace{11mm}
{\Large\color{accent}Силы, трение и второй закон Ньютона}\par
\vspace{8mm}
{\normalsize Горизонтальное движение бруска}\par
\vfill
{\large Володя Хачатурян}\par
\vspace{11mm}
16 ноября 2026 г.\par
\vspace{25mm}
{\small Лёвин Артём Александрович}\par
{\small эксклюзивно для Володи Хачатуряна}\par
\vspace{24mm}
\begin{tikzpicture}[x=1cm,y=1cm]
\draw[accent,line width=0.85pt] (0,0) .. controls (1.1,0.43) and (2.0,-0.26) .. (3.0,0)
  .. controls (4.2,0.35) and (5.0,-0.32) .. (6.1,0);
\end{tikzpicture}
\vspace*{8mm}
\end{center}
\clearpage

{\Large\bfseries\color{deepaccent}Карта разбора}\par
\vspace{1mm}
В таблицу включены задания, где решение отсутствует, не завершено или содержит подтверждённую ошибку в рассуждении. Для каждого такого задания приведено исправление и одна новая задача на закрепление.

\renewcommand{\arraystretch}{1.25}
\small
\begin{tabularx}{\linewidth}{@{}>{\centering\arraybackslash}p{11mm}p{25mm}X>{\raggedright\arraybackslash}p{25mm}>{\raggedright\arraybackslash}p{29mm}@{}}
\toprule
№ & Тема & Суть проблемы & Ваш ответ & Правильный результат \\
\midrule
\hyperref[task:1]{1} & Проекции силы & Решение отсутствует; для чисел не указан угол. & Не указан & $F_x=16\cos\alpha$, $F_y=16\sin\alpha$ при силе вверх-вправо \\
\hyperref[task:2]{2} & Сила опоры & Неверный знак уравнения; не учтён наклон тяги. & $30\,\text{Н}$ & Нужен угол наклона \\
\hyperref[task:7]{7} & Баланс сил & Нет расчёта сил и ускорения. & Не указан & $a_x=0$ \\
\hyperref[task:9]{9} & Обратная задача & Не найдены сила трения и сила тяги. & Не указан & $F=8\,\text{Н}$ \\
\hyperref[task:10]{10} & Знак ускорения & Не выполнена проекция второго закона Ньютона. & Не указан & $a_x=-3{,}5\,\text{м/с}^2$ \\
\bottomrule
\end{tabularx}
\normalsize

\vspace{5mm}
\subblock{Условные обозначения}
Во всех вычислениях ниже используется $g=10\,\text{м/с}^2$ --- то же значение, которое принято в представленных записях. Ось $Ox$ направляем вправо, ось $Oy$ --- вверх. $N$ обозначает модуль реакции опоры; $F_{\text{тр}}$ --- модуль силы трения скольжения.

\clearpage
\thispagestyle{empty}
\vspace*{0.34\textheight}
\begin{center}
{\fontsize{26}{32}\selectfont\color{deepaccent}Разбор ошибок}\par
\vspace{7mm}
{\normalsize Последний верный шаг, точка затруднения, правило, исправление и самостоятельная тренировка}
\end{center}
\clearpage

\phantomsection\label{task:1}
\reviewtitle{1}
\subblock{Текст задания}
На брусок действует наклонная сила $F=16\,\text{Н}$. Найдите её проекции $F_x$ и $F_y$.

\subblock{Ваше решение}
Записано данное $F=16\,\text{Н}$ и обозначены искомые проекции $F_x$ и $F_y$. Расчёт и ответ отсутствуют.

\subblock{Ваш ответ}
Не указан.

\subblock{Ошибка}
\errorlabel\ Решение остановлено до разложения вектора силы на проекции. Кроме того, исходное условие не задаёт угол наклона и ориентацию силы относительно осей, поэтому единственного численного ответа получить нельзя.

\subblock{Где именно возникает затруднение}
Последний правильный шаг --- выделение известного модуля $F=16\,\text{Н}$ и неизвестных $F_x$, $F_y$. Первый отсутствующий шаг --- выбор осей и обозначение угла между вектором силы и положительным направлением оси $Ox$.

\subblock{Почему без этого шага нельзя получить результат}
Проекция описывает, какая часть вектора направлена вдоль конкретной оси. При фиксированном модуле $16\,\text{Н}$ горизонтальная и вертикальная части меняются с углом наклона. Например, у горизонтального вектора вертикальная проекция равна нулю; при наклоне вверх обе проекции могут быть положительными. Знания одного модуля силы для двух численных проекций недостаточно.

\subblock{Нарушенное правило}
Проекции вектора определяются его модулем, направлением и выбранными осями. Знак каждой проекции зависит от направления вектора относительно положительного направления соответствующей оси.

\subblock{Ключевая идея}
\keyidea\ Если угол $\alpha$ отсчитывается от положительного направления оси $Ox$ к силе, направленной вправо-вверх, то прилежащий катет соответствует горизонтальной проекции, а противолежащий --- вертикальной:
\[
\cos\alpha=\frac{F_x}{F},\qquad
\sin\alpha=\frac{F_y}{F}.
\]

\subblock{Исправление от последнего правильного шага}
К записи $F=16\,\text{Н}$ добавьте угол $\alpha$ и договорённость об осях. Тогда обе проекции выражаются через заданный модуль силы и тригонометрические функции этого угла.

\subblock{Полное правильное решение}
\begin{enumerate}
\item Направим $Ox$ вправо, $Oy$ вверх.
\item Пусть вектор силы направлен вправо-вверх под углом $\alpha$ к горизонтали.
\item По определению косинуса в прямоугольном треугольнике $F_x=F\cos\alpha$.
\item По определению синуса $F_y=F\sin\alpha$.
\item Подставим модуль силы:
\[
\boxed{F_x=16\cos\alpha\,\text{Н}},\qquad
\boxed{F_y=16\sin\alpha\,\text{Н}}.
\]
\end{enumerate}
Если сила направлена вправо-вниз, её вертикальная проекция отрицательна: $F_y=-16\sin\alpha$ при $\alpha$ как остром угле к горизонтали. Аналогично изменяется знак $F_x$ при направлении влево.

\subblock{Проверка}
Для любого угла сумма квадратов перпендикулярных проекций должна равняться квадрату модуля:
\[
F_x^2+F_y^2=16^2(\cos^2\alpha+\sin^2\alpha)=256.
\]
При $\alpha=0^\circ$ получается $F_x=16\,\text{Н}$, $F_y=0$, что соответствует горизонтальной силе.

\subblock{Итог}
\resultlabel\ Формулы для проекций получены. Для численного ответа необходимо дополнить условие углом и направлением силы.

\subblock{Задача на закрепление}
Сила $20\,\text{Н}$ направлена вправо-вверх под углом $30^\circ$ к горизонтали. Ось $Ox$ направлена вправо, ось $Oy$ вверх. Найдите $F_x$ и $F_y$.

\clearpage
\phantomsection\label{task:2}
\reviewtitle{2}
\subblock{Текст задания}
Масса бруска $3\,\text{кг}$, модуль наклонной тяги $10\,\text{Н}$. Найдите реакцию опоры.

\subblock{Ваше решение}
Записано $N-mg>0$, затем сделан переход к $N=mg=30\,\text{Н}$. В полученном ответе используется $g=10\,\text{м/с}^2$.

\subblock{Ваш ответ}
$N=30\,\text{Н}$.

\subblock{Ошибка}
\errorlabel\ В первой записи использован знак строгого неравенства между $N-mg$ и нулём. Из такого неравенства нельзя получить равенство $N=mg$. Численный ответ $30\,\text{Н}$ при горизонтальной тяге верен, но запись его обоснования содержит ошибку. Дополнительно слово «наклонная» в условии требует учитывать вертикальную проекцию тяги, если она отлична от нуля.

\subblock{Где именно возникает ошибка}
Последний правильный шаг --- выделение вертикальных сил: реакции опоры и тяжести. Первый неверный шаг --- $N-mg>0$. При отсутствии вертикального ускорения должно выполняться равенство суммы проекций нулю. Для наклонной тяги дополнительно появляется $F_y$.

\subblock{Почему записанный переход неверен}
Неравенство $N-mg>0$ означает $N>mg$, то есть реакция опоры больше силы тяжести. Следующий переход $N=mg$ одновременно с этим неравенством несовместим. Кроме того, наклонная сила может частично разгружать опору при направлении вверх. Без величины угла и направления тяги нельзя определить $F_y$ и окончательный численный ответ для действительно наклонной силы.

\subblock{Нарушенное правило}
Второй закон Ньютона записывают для каждой оси отдельно. При $a_y=0$ сумма проекций сил на $Oy$ равна нулю:
\[\sum F_y=0.\]

\subblock{Ключевая идея}
\keyidea\ Если тяга горизонтальна, $F_y=0$ и $N=mg$. Если тяга наклонена вверх, $N+F_y-mg=0$, поэтому $N=mg-F_y$. Реакция опоры зависит от вертикальной составляющей тяги.

\subblock{Исправление от последнего правильного шага}
Вместо $N-mg>0$ запишите уравнение по вертикали. При горизонтальной тяге:
\[N-mg=0.\]
При наклонной тяге вверх:
\[N+F\sin\alpha-mg=0,\]
если $\alpha$ --- угол к горизонтали.

\subblock{Полное правильное решение}
\textbf{Случай горизонтальной тяги.} Брусок остаётся на горизонтальной поверхности, поэтому $a_y=0$. Силы по вертикали: $N$ вверх и $mg$ вниз. Следовательно,
\[
N-mg=0,\qquad N=mg=3\cdot10=\boxed{30\,\text{Н}}.
\]
\textbf{Случай наклонной тяги вверх.} Появляется положительная вертикальная проекция $F\sin\alpha$. Тогда
\[
N+F\sin\alpha-mg=0,\qquad
\boxed{N=30-10\sin\alpha\,\text{Н}}.
\]
Угол в исходном тексте не дан; направление наклонной силы также не уточнено. Следовательно, численное значение реакции опоры в таком случае остаётся неопределённым. Для тяги вниз при том же модуле вертикальная проекция отрицательна, и реакция опоры увеличивается.

\subblock{Проверка}
При $\alpha=0^\circ$ формула наклонной тяги вверх даёт $N=30\,\text{Н}$. При $\alpha=30^\circ$ получится $N=30-5=25\,\text{Н}$, что согласуется с частичным подъёмом бруска силой тяги.

\subblock{Итог}
\resultlabel\ Уравнение по вертикали должно содержать знак равенства. Ответ $30\,\text{Н}$ относится к случаю нулевой вертикальной проекции тяги; при наклонной тяге требуется знать угол и направление.

\subblock{Задача на закрепление}
Брусок массой $4\,\text{кг}$ находится на горизонтальной опоре. К нему приложена сила $12\,\text{Н}$ под углом $30^\circ$ к горизонтали, направленная вправо-вверх. Найдите реакцию опоры при $g=10\,\text{м/с}^2$.

\clearpage
\phantomsection\label{task:7}
\reviewtitle{7}
\subblock{Текст задания}
Брусок массой $3\,\text{кг}$ скользит вправо. Сила тяги $12\,\text{Н}$, коэффициент трения скольжения $\mu=0{,}4$. Найдите проекцию ускорения $a_x$.

\subblock{Ваше решение}
Под заголовком задания оставлен знак вопроса. Расчёт сил, уравнение Ньютона и ответ отсутствуют.

\subblock{Ваш ответ}
Не указан.

\subblock{Ошибка}
\errorlabel\ Решение не завершено: отсутствует переход от данных к горизонтальной равнодействующей сил.

\subblock{Где именно возникает затруднение}
Последний правильный шаг --- распознавание номера задания и постановка задачи. Первый отсутствующий шаг --- расчёт реакции опоры и силы трения скольжения.

\subblock{Почему этот шаг необходим}
Брусок движется по горизонтали под действием тяги вправо и трения влево. Ускорение определяется разностью этих сил. Нельзя по одной лишь величине тяги заключить, что ускорение положительно: трение может скомпенсировать тягу полностью.

\subblock{Нарушенное правило}
Нужно использовать второй закон Ньютона в проекциях: $\sum F_x=ma_x$. Сила трения скольжения по модулю равна $F_{\text{тр}}=\mu N$, а для горизонтальной поверхности без вертикальной составляющей тяги $N=mg$.

\subblock{Ключевая идея}
\keyidea\ Сначала найти трение, затем равнодействующую, затем разделить её на массу. Если тяга и трение равны по модулю, горизонтальное ускорение равно нулю.

\subblock{Исправление от последнего правильного шага}
К данным $m=3\,\text{кг}$, $F=12\,\text{Н}$, $\mu=0{,}4$ добавьте расчёт $N=mg$, затем $F_{\text{тр}}=\mu N$. Сразу после этого запишите $F-F_{\text{тр}}=ma_x$ с учётом направлений сил.

\subblock{Полное правильное решение}
\begin{enumerate}
\item Направим $Ox$ вправо. Брусок скользит вправо, поэтому трение скольжения направлено влево.
\item По вертикали ускорение равно нулю; при горизонтальной тяге:
\[N-mg=0,\qquad N=3\cdot10=30\,\text{Н}.\]
\item Найдём модуль силы трения:
\[F_{\text{тр}}=\mu N=0{,}4\cdot30=12\,\text{Н}.\]
\item Запишем второй закон Ньютона по горизонтали:
\[F-F_{\text{тр}}=ma_x.\]
\item Подставим числа и решим уравнение:
\[12-12=3a_x,\qquad 0=3a_x,\qquad\boxed{a_x=0}.\]
\end{enumerate}

\subblock{Проверка}
Сумма проекций горизонтальных сил равна $12-12=0\,\text{Н}$. При массе $3\,\text{кг}$ из второго закона Ньютона следует $a_x=0$. Это означает постоянную скорость в рассматриваемом промежутке времени; брусок может продолжать скользить вправо.

\subblock{Итог}
\resultlabel\ $\boxed{a_x=0\,\text{м/с}^2}$. Направленное вправо движение совместимо с нулевым ускорением.

\subblock{Задача на закрепление}
Брусок массой $2\,\text{кг}$ скользит вправо по горизонтальной поверхности. На него действует горизонтальная сила $6\,\text{Н}$ вправо; коэффициент трения скольжения равен $0{,}3$. Найдите $a_x$ при $g=10\,\text{м/с}^2$.

\clearpage
\phantomsection\label{task:9}
\reviewtitle{9}
\subblock{Текст задания}
Брусок массой $2\,\text{кг}$ скользит вправо с ускорением $2\,\text{м/с}^2$. Коэффициент трения $\mu=0{,}2$. Найдите модуль горизонтальной силы тяги $F$.

\subblock{Ваше решение}
Выписаны исходные величины: $m=2\,\text{кг}$, $a_x=2\,\text{м/с}^2$, $\mu=0{,}2$. Далее поставлен знак вопроса. Уравнение и численный ответ отсутствуют.

\subblock{Ваш ответ}
Не указан.

\subblock{Ошибка}
\errorlabel\ Решение остановлено перед применением второго закона Ньютона к неизвестной силе тяги.

\subblock{Где именно возникает затруднение}
Последний правильный шаг --- запись массы, ускорения и коэффициента трения. Первый отсутствующий шаг --- расчёт $N=mg$ и $F_{\text{тр}}=\mu N$, после которого можно выразить искомую силу из $F-F_{\text{тр}}=ma_x$.

\subblock{Почему необходим этот шаг}
На создание ускорения действует равнодействующая сила. Часть тяги уравновешивает трение, а оставшаяся часть разгоняет брусок. Поэтому модуль тяги должен превышать произведение $ma_x$, когда ускорение направлено вправо и трение действует влево.

\subblock{Нарушенное правило}
При ускорении вправо и движении вправо второй закон Ньютона в проекции на $Ox$ имеет вид $F-F_{\text{тр}}=ma_x$. Отсюда $F=ma_x+F_{\text{тр}}$.

\subblock{Ключевая идея}
\keyidea\ Это обратная задача: ускорение уже задано, поэтому нужно найти сумму силы, которая создаёт ускорение, и силы, преодолевающей трение.

\subblock{Исправление от последнего правильного шага}
Сохраните записанные данные. Добавьте $g=10\,\text{м/с}^2$, найдите реакцию опоры и трение, затем выразите $F$ из проекции второго закона Ньютона.

\subblock{Полное правильное решение}
\begin{enumerate}
\item Направим $Ox$ вправо. Тяга действует вправо, трение скольжения --- влево.
\item Так как вертикального ускорения нет, $N=mg=2\cdot10=20\,\text{Н}$.
\item Найдём силу трения скольжения:
\[F_{\text{тр}}=\mu N=0{,}2\cdot20=4\,\text{Н}.\]
\item Второй закон Ньютона по оси $Ox$:
\[F-F_{\text{тр}}=ma_x.\]
\item Выразим искомую силу и подставим значения:
\[
F=ma_x+F_{\text{тр}}=2\cdot2+4=\boxed{8\,\text{Н}}.
\]
\end{enumerate}

\subblock{Проверка}
Возьмём найденную тягу $8\,\text{Н}$ и вычтем трение $4\,\text{Н}$. Равнодействующая равна $4\,\text{Н}$. Тогда $a_x=4/2=2\,\text{м/с}^2$, как в условии.

\subblock{Итог}
\resultlabel\ $\boxed{F=8\,\text{Н}}$. Тяга расходуется на компенсацию трения и создание положительного ускорения.

\subblock{Задача на закрепление}
Брусок массой $3\,\text{кг}$ движется вправо с ускорением $1{,}5\,\text{м/с}^2$ под действием горизонтальной силы вправо. Коэффициент трения скольжения равен $0{,}2$. Найдите модуль силы тяги при $g=10\,\text{м/с}^2$.

\clearpage
\phantomsection\label{task:10}
\reviewtitle{10}
\subblock{Текст задания}
Брусок массой $4\,\text{кг}$ скользит вправо. Сила тяги $10\,\text{Н}$, коэффициент трения скольжения $\mu=0{,}6$. Найдите $a_x$ и объясните физический смысл знака.

\subblock{Ваше решение}
Под номером задания поставлен знак вопроса. Силы, ускорение и ответ не вычислены.

\subblock{Ваш ответ}
Не указан.

\subblock{Ошибка}
\errorlabel\ Не выполнено сравнение тяги с трением и не определён знак проекции ускорения.

\subblock{Где именно возникает затруднение}
Последний правильный шаг --- указание номера задания. Первый отсутствующий шаг --- нахождение $N$ и силы трения, позволяющее установить, какая сила больше.

\subblock{Почему этот шаг необходим}
Условие сообщает, что брусок скользит вправо. Это описывает направление скорости в рассматриваемый момент. Ускорение характеризует изменение скорости и может быть направлено против движения. Если трение влево превосходит тягу вправо, равнодействующая и ускорение направлены влево.

\subblock{Нарушенное правило}
Знак ускорения определяется знаком суммы проекций сил: $ma_x=F-F_{\text{тр}}$. Направление скорости само по себе знак ускорения не фиксирует.

\subblock{Ключевая идея}
\keyidea\ Сравните $F$ и $\mu mg$. Если тяга меньше трения, то $F-F_{\text{тр}}<0$ и $a_x<0$ при оси вправо. Такое ускорение уменьшает положительную скорость.

\subblock{Исправление от последнего правильного шага}
После данных $m=4\,\text{кг}$, $F=10\,\text{Н}$, $\mu=0{,}6$ найдите $N=mg=40\,\text{Н}$ и $F_{\text{тр}}=24\,\text{Н}$. Затем подставьте обе силы в закон Ньютона с правильными знаками.

\subblock{Полное правильное решение}
\begin{enumerate}
\item Ось $Ox$ направим вправо. Брусок скользит вправо, поэтому сила трения скольжения направлена влево.
\item По вертикали ускорения нет, поэтому $N=mg=4\cdot10=40\,\text{Н}$.
\item Найдём модуль силы трения:
\[F_{\text{тр}}=\mu N=0{,}6\cdot40=24\,\text{Н}.\]
\item Тяга вправо даёт положительную проекцию $+10\,\text{Н}$, а трение влево --- отрицательную $-24\,\text{Н}$.
\item Запишем второй закон Ньютона и вычислим ускорение:
\[
10-24=4a_x,\qquad -14=4a_x,\qquad
\boxed{a_x=-3{,}5\,\text{м/с}^2}.
\]
\end{enumerate}

\subblock{Проверка}
Равнодействующая горизонтальных сил $-14\,\text{Н}$ направлена влево; после деления на положительную массу $4\,\text{кг}$ знак ускорения сохраняется отрицательным. При начальной скорости вправо $v_{0x}>0$ на участке продолжающегося скольжения имеем $v_x(t)=v_{0x}-3{,}5t$: скорость уменьшается до остановки. Поведение после остановки требует отдельного учёта трения покоя; данных для него здесь нет.

\subblock{Итог}
\resultlabel\ $\boxed{a_x=-3{,}5\,\text{м/с}^2}$. Отрицательный знак означает ускорение влево при выбранной оси $Ox$; брусок в этот момент замедляется, продолжая скользить вправо.

\subblock{Задача на закрепление}
Брусок массой $5\,\text{кг}$ скользит вправо по горизонтальной поверхности. На него действует горизонтальная сила $12\,\text{Н}$ вправо, коэффициент трения скольжения равен $0{,}4$. Найдите $a_x$ и объясните знак ускорения при $g=10\,\text{м/с}^2$.

\end{document}